\subsection{二積追加型の片外側Ore表示と反交換二次元}
\label{subsec:cubic-CY-one-segre-double-split}

\begin{proposition}\label{prop:cubic-CY-one-segre-double-split-outer-center}
一交換子の非共有Segre射影 $ab+cd$ について、固定出典\cite{AS4CYpinned}の
非等方片外側三族 $C,N,Z$ を考える。
整域切断を許し得るのは $Z(s,0)$ の $s^2=-1/12$ のみである。
このとき $q=(2s+1)/(2s-1)$ は原始三乗根で、
全中央三次は
\[
 G=P(a,b)+C c^3+D d^3,\qquad P\in k[a,b]_3
\]
である。また同出典の対称外側・非分解四交差族には中央三次がなく、
H3包絡は存在しない。
\end{proposition}
\begin{proof}
非等方片外側のOre表示を用いる。$X=(a,b)^{\mathsf T}$、
$\kappa=s^2+4(pr-h^2)+1/4$ とおく。
$cX=MXc$、$dX=M^{-1}Xd$、$dc=\delta_0 cd+Kab$ であり、
\[
 M=\frac{1}{\kappa-s}
 \begin{pmatrix}\kappa-2h-1/2&-2r\\2p&\kappa+2h-1/2\end{pmatrix},
 \qquad
 \delta_0=\det M=\frac{\kappa+s}{\kappa-s},\quad
 K=\frac{2s}{\kappa-s}.
\]
$s\ne0$、$\kappa^2\ne s^2$、$\delta_0\ne1$ であり、
$\operatorname{tr}M-(1+\delta_0)=-1/(\kappa-s)$。
従って $M$ は固有値1や $\delta_0$ を持たない。
$M^2=I$ なら $\delta_0=-1$、$\kappa=0$ だが、
このとき $\operatorname{tr}M=1/s\ne0$ で矛盾する。

$c^i d^j$ の順序基底により、$a,b$ と可換な三次元の各項は
$M^{i-j}=I$ を満たす。
$M^3\ne I$ なら候補は $P+Qcd$（$\deg Q=1$）である。
$c$ との可換性の $cd$ 係数は $\sigma_M(Q)=\delta_0 Q$ であり、
$\delta_0$ が固有値でないため $Q=0$。
$P(a,b)$ は可換二元三次なので因数分解して整域切断を与えない。
$M^3=I$ なら固有値は三乗根で、$\delta_0\ne1$ かつ固有値1を持たないため
$M=qI$ となる。従って $p=r=h=0$、$\kappa=s^2+1/4$ で、
$q=(2s+1)/(2s-1)$、$q^2+q+1=0$ は $s^2=-1/12$ と同値である。
この値では $c^3,d^3$ が中央になる。
例えば $dc^3$ をOre式で展開すると、補正係数は
$K(1+\delta_0^2+\delta_0^{-2})=0$ である。$d^3c$ も交換して得られる。
従って残る $Qcd$ は先と同様に消え、$P,c^3,d^3$ の全ペンシルを得る。

対称外側・非分解四交差族では、固定出典の (FSC8)--(FSC9) が
全一次元と反交換する非零二次元 $Q_0$ を与える。
$G$ が中央三次なら $Q_0G=-GQ_0= -Q_0G$ なので、
包絡の整域性から $G=0$ となる。
\end{proof}

\begin{theorem}\label{thm:cubic-CY-one-segre-double-split-outer-H3}
前命題の $Z(s,0)$ の共鳴点で、$\nu=qK/(1-q)$ とする。
切断 $P+C c^3+Dd^3$ がH3となる必要十分条件は次の通りである。
\begin{enumerate}
\item $CD\ne0$。
\item $P$ は $[1:0],[0:1]\in\mathbb P^1$ で消える場合、そこで単純に消える。
\item $\Delta(z)=P(1,z)^2-4CD\nu^3 z^3$ は $k^*$ 上に重根を持たない。
\end{enumerate}
二つの $s$ は $c,d$ の交換で同型である。
一方に固定して $C=D=1$ と規格化した幾何学的群oidは、
上の $P$ の集合に
\[
 \{(u,v)\in(k^*)^2:(uv)^3=1\}\rtimes C_2
\]
が対角変換と $a,b$ の交換で作用する群oidである。
一般慣性は $\mu_3$、次元は3である。例えば $P=a^3+2b^3$、$C=D=1$ はH3を与える。
\end{theorem}
\begin{proof}
$C=0$ なら $c$ 軸はshift固定点対象であり、切断はその一次イデアルの平方に入る。
$c$ の作用を恒等として一次化すると $a,b,d$ の三方向が自由で、
次次数への倍率は $q,q,\delta_0^{-1}=q$ である。
同様に $D=0$ の $d$ 軸では三倍率が $q^{-1}$。
行列式は1なので奇数Serre障害に反する。

$CD\ne0$ とする。$a,b$ を固定し $c,d$ を $q,q^{-1}$ 倍する自己同型で
右Zhangねじりを取ると、$a,b$ は中央となり
\[
 dc=qcd+(1-q)\nu ab .
\]
三次切断の式は変わらない。$X=c^3,Y=d^3,w=cd-\nu ab$ と書くと
$wc=qcw$、$w^3=XY-\nu^3 a^3b^3$ である。
切断 $R'$ は $Y=-(P+CX)/D$ により中心 $k[a,b,X]$ 上で階数9の自由加群となる。

$a$ 軸で $P(1,b)$ が $b^2$ で割れるとする。
$a=1$ の次数零環を閉点 $b=X=Y=0$ で完備化する。
$c,d$ の完備自由環の中で
$b=(dc-qcd)/((1-q)\nu)$ を消去すると、
$b$ の二中央性関係と $P(1,b)+C c^3+D d^3$ が表示関係になる。
この完備化と $(c,d)$ による完備化は一致する。
実際、中央極大イデアルの剰余環では $b=0,c^3=d^3=0$ で
量子平面の有限次元局所商を得るので、そのradicalは冪零である。
三つの関係は全て最低次数3であり、
それぞれの最低記号の二重次数は $(2,1),(1,2)$ と純三乗の部分なので独立である。
従って極小関係空間の次元は3以上であり、
剰余単純の $\operatorname{Ext}^2$ の次元も3以上となる。
これは増大イデアルを生成する二元で自由解消を始めれば、
二番目の極小生成元がちょうど極小関係に対応することから従う。
中心有限性により閉点完備化はこの単純のExtを変えず、
$a$ の局所化もその点のtails Extを変えない。
後者は加群 $k[a]$ を $k[a^{\pm1}]$ へ送るtorsion商と
Ore局所化の導来随伴で従う。
ねじり前では同じ点がshiftで固定され、H3のSerre双対による
$\dim \operatorname{Ext}^2=1$ に反する。
$b$ 軸も同様である。

$c$ のchartは定理\ref{thm:cubic-CY-one-segre-square-split-resonance-exact}と
同じ巡回order模型で、今度は
\[
 F(\xi,\eta)=-(C+P(\xi,\eta))/D-\nu^3\xi^3\eta^3 .
\]
$d$ のchartも同様である。
$F=0$ は原点を通らず、座標軸上での零があれば、
一変数の非零定数と三乗の和なので単純である。
トーラス上では $z=\eta/\xi,x=\xi^{-3}$ を用いて、
$F=0$ は
\[
 g(x,z)=C x^2+P(1,z)x+D\nu^3 z^3=0
\]
のエタール引戻しとなる。
$g=g_x=g_z=0$ は $\Delta$ のトーラス上の重根と同値である。
同じ循環交差積の局所計算で、これが有限大域次元に対する
正確な障害となる。

三条件から整域性も従う。$a=1$ の中心分数体上での巡回代数は
$x$ と $-g/D$ をパラメータに持つ。
$P(1,0)$ が零なら単純であり、$z=0$ 上の $g=0$ の二枝で
$x$ の位数は1と2である。実際、$x=zX_1$ とおくと
$CX_1^2+(P(1,z)/z)X_1+D\nu^3 z=0$ の特殊繊維は二単根を持つ。
$P(0,1)$ のみが零の場合も、$a,b$ を交換して同じ議論を適用する。
従って該当成分の関数体上で $x$ は立方でなく、
$g$ の既約因子は重複しないため、その因子の付値で $-g/D$ の位数は1である。
両座標軸で $P$ が非零なら $\Delta$ は次数6の平方自由多項式であり、
$g=0$ の滑らかな射影モデルは種数2となる。
$x$ の唯一の極は無限遠の一枝で位数3であり、
有限点では極を持たない。$x$ が立方なら立方根は次数1の写像を与え、
種数2に反する。
純超越変数 $a$ を加えても剰余関数体での非立方性は保たれる。
いずれもノルムの付値により巡回代数は非分裂、従って次数3の斜体である。
元の斉次環は中心上自由なので斜体に入る。
ねじりも整域性を保つ。

残る $X=Y=0$ のうち $ab\ne0$ では、
$w^3=-\nu^3a^3b^3$ の三冪等元が分裂し、
$c,d$ の二つの単元辺が全冪等元を同値にする。
自由階数9から完備化は中心局所環上の行列環である。
座標軸の点では単純零条件により、もう一方の中央基底変数を
$X,Y$ から陰関数で消去できる。
例えば $a=1$ なら $b\in(X,Y)k[[X,Y]]$ であり、
量子平面の降下重みでOre補正の重みは3以上。
完成した環の随伴環は量子平面なので、左右大域次元は高々2である。
こうして全中心chartで有限大域次元を得る。
中心上有限な整域のrank-one、局所Extの貼合せ、AS Gorenstein双対と
Zhang同値から、残るH3条件は平方・積追加の場合と同様に従う。

最後に固定出典の (DN3)--(DN4) より、固定した $s$ の全包絡自己同型は
$a,b$ の対角変換と交換、および $c\mapsto zc,d\mapsto uv,d/z$ である。
非零な二純三乗係数はこれらにより1へ規格化できる。
共通次数スカラーで $z=1$ とすれば、固定条件は $(uv)^3=1$ である。
従って全矢は表示の群で尽くされる。
一般の四係数が非零の $P$ では対角安定化に $u=v\in\mu_3$ が必要である。
交換付きの固定条件は四係数を順に $A_i$ と書くと
$A_0 A_2^3=A_3 A_1^3$ を含むので、一般点では成立しない。
係数空間の次元4から群の次元1を引いて次元3である。
$\nu^3=1$ は有限数体の恒等式なので、表示した例の判別式は $1+4z^6$ になる。
\end{proof}

Ore trace、原始三乗根条件、巡回ノルムおよび軸での三極小関係の独立性は
\par\noindent\texttt{verify\_cubic\_CY\_one\_segre\_double\_split\_outer.py}
\par
で検証する。

\begin{proposition}\label{prop:cubic-CY-one-segre-double-split-isolated-empty}
固定出典\cite{AS4CYpinned}の二積追加射影にある五つの孤立包絡には、
H3切断がない。
\end{proposition}
\begin{proof}
孤立五表示は両外側 (DE2)、一交差 (DC1)、
対角二交差の例外体積 (DX1) と八分の一体積の二表示 (EN3) である。
元の六関係から次数2,3,4の商次元 $(10,20,35)$ を確認し、
中心条件 $\mathcal B_3\to\mathcal B_4^{\oplus4}$ の階数を厳密に求めると、
(DX1) は18、他の四表示は20である。

(DX1) で $\gamma_0=i\sqrt3$、$\zeta=(-1+\gamma_0)/2$ とおき、
\[
 \ell=a+\gamma_0d/18,\qquad m=b-6\gamma_0 c
\]
を取る。元の二次商で $m\ell=\zeta^2\ell m$ が成立する。
二つの立方 $\ell^3,m^3$ は独立かつ中央であることを
次数3と4の商で確認する。
中心次元は2なので、この二元が全中央三次を張る。
$\zeta^2$ は原始三乗根であるから
$(u\ell+v m)^3=u^3\ell^3+v^3m^3$。
代数閉性により任意の中央三次はこの形となり、
非零な一次因子は三次切断で消えないため整域切断とならない。
全計算を公開検証器
\par\noindent\texttt{verify\_cubic\_CY\_one\_segre\_double\_split\_isolated.py}
\par\noindent
で、ガウス体と $\mathbb Q(i\sqrt3)$ 上で再現できる。
\end{proof}

\begin{proposition}\label{prop:cubic-CY-one-segre-double-split-small-empty}
固定出典\cite{AS4CYpinned}の二積追加型で、
隣接二交差 (DA1)、対角二交差の三連続族 (DD1),(EN2),(EB2)、
非分解三交差 (TN1)、外側零の四交差 (FZ1) と、等方片外側の九族にはH3切断がない。
\end{proposition}
\begin{proof}
最初の五族について元の六関係から次数2,3,4の商と
中心行列を作る。商座標写像の全分母と、中心行列の選んだ
階数20の小行列式の分子・分母の因子を記録する。
(DA1),(DD1),(EB2) では全て元の許容域の単元である。
(EN2) では七点、(TN1) では二点が残るが、
元の六関係から再計算して全て階数20となる。
(FZ1) では $y-z\ne0$ の余分な因子は、
$y-z=\pm1/512$ と
\[
 (y-z)(y+3z)=1/262144,\qquad (y-z)(y+z)=1/262144
\]
である。前二つは直線、後二つは $y-z$ を座標にして有理化できる。
その四表示の全例外値も再計算して階数20となる。

等方片外側では固定出典の二行列 $H_a,H_n$ を一括して扱う。
$H_a$ では $e=a$、$H_n$ では $e=a+b$ と置き、
生成元 $e,b,c,d$ の順序基底を使う。$q=(4s+1)/(4s-1)$ とすれば
$q\ne0,\pm1$ であり、$e,b$ は可換、$e$ は $c,d$ と反交換する。
元関係を同じ順序で書き直すと
\begin{align*}
 cb&=qbc+e(uc+vd),&
 db&=q^{-1}bd+e(wc+xd),\\
 dc&=-qcd+K_a eb+K_b b^2 .
\end{align*}
ここで $u,v,w,x,K_a,K_b$ はスカラーである。
$v,w$ の零化を含め、この表示は二つの元関係行列の基底の取り替えで得られる。
基底変更の行列式は $(4s-1)^2(4s+1)/(128s)$ で、
全許容値で可逆である。

$e$ と可換な三次式は $c,d$ の合計次数が偶数の項だけを持つ。
$e$ を法とした商で $b$ は $c,d$ をそれぞれ $q,q^{-1}$ 倍するので、
$b$ と可換な混合項としては $bcd$ しか残らない。
$c$ との交換子で $bc^2d$ の係数を比較すると
$q$ と $-q$ の差によりこれも零となる。
従って元の中央三次は
$P(e,b)+e(\alpha c^2+\beta cd+\gamma d^2)$ と書ける。
$b$ との交換子を $e^2$ を法として比較すると、
$q^2\ne1$ から $\alpha=\gamma=0$。
最後に $c$ との交換子の $ec^2d$ 係数は、
$ce=-ec$ と $dc=-qcd+K_a eb+K_b b^2$ から $(q-1)\beta$ になる。
$P$ の交換子は $c,d$ の合計次数が1以下なのでこの係数に寄与しない。
よって $\beta=0$ であり、全中央三次は可換二元三次 $P(e,b)$ となる。
代数閉体上で因数分解するので整域切断を与えない。

中心の小行列式の例外曲線・例外点と、等方片外側の元関係の書換えを
公開検証器で再現する。
\par\noindent\texttt{verify\_cubic\_CY\_one\_segre\_double\_split\_remaining\_small.py}
\par
\end{proof}

\begin{proposition}\label{prop:cubic-CY-one-segre-double-split-quadratic-reduction}
残る分解三交差 (TD1) と片外側四交差 (FO1)--(FO2) において、
実交換子平面から非零中央二次 $Q_0$ が得られる。
次の記法では $U,V$ は可換で、$Q_0$ の係数は三つの表示座標の一つに依存しない。
\begin{enumerate}
\item (TD1) の $(s,u,w)=(p,u,w)$、$t=p-1/4$ では
\[
 U=a+tc,\qquad V=pb+c+d,\qquad
 Q_0=p^2U^2+pw(16p+64w-2)UV+w^2V^2.
\]
\item (FO1)--(FO2) の $(t,\sigma,q)=(p,u,w)$、
$h=1/64-p^2/4$、$\eta=pw+1/32$ では
\begin{align*}
 U&=a+hb+p(c-d)/2,\qquad V=c+d,\\
 Q_0&=pU^2+2\eta UV+16p\left(\eta^2+(1-16p^2)/1024\right)V^2.
\end{align*}
\end{enumerate}
いずれも $Z(\mathcal B)_3\to Z(\mathcal B/(Q_0))_3$ は単射である。
二次商の次数2,3,4の次元は $(9,16,25)$ である。
\end{proposition}
\begin{proof}
$[U,V]$ は元の実交換子そのものである。
公開検証器は元の六関係から次数2,3の商を構成し、
全一般係数で $[Q_0,a]=[Q_0,b]=[Q_0,c]=[Q_0,d]=0$
を関数体上で再生する。
元の関係の係数は固定出典の許容開集合上で正則であり、
次数三商はAS4のHilbert値によりその全域で一定次元である。
したがって上の交換子は全許容値で零である。
$U,V$ は独立なので、整域である元の包絡内でその可換二元二次式の
各一次因子は非零、従って $Q_0\ne0$ である。

一次中心があれば、非零一次元 $\ell$ に対し
$\ell\wedge \mathcal B_1$ の三次元が二次関係の交代部分に入るが、
交代部分は一次元なので矛盾する。
中央三次が二次商で消えれば $G=Q_0\ell$ であり、
$Q_0$ の正則性から $\ell$ も中央となるため $G=0$ である。
二次商のHilbert級数は $(1-z^2)/(1-z)^4$ である。
\end{proof}

公開検証器は次である。
\par\noindent\texttt{verify\_cubic\_CY\_one\_segre\_double\_split\_quadratics.py}
\par
この還元は全パラメータで使える。二次商の中心三次と元の包絡への
持ち上げの全判定は引き続き必要である。

(TD1) と (FO1)--(FO2) の関数体上では、元の六関係から作った次数4商で
中心行列の階数20を確認した。
用いた小行列式の分子・分母および商の射影分母の全既約因子を
検証器の \texttt{--cubic TD1} または \texttt{--cubic FO1} で再生する。
従ってこの係数域で未判定の中央三次は、証明書にある因子の零集合に限られる。
この計算はそれらの零集合を判定するものではない。
